\documentclass[12pt, a4paper]{article}

% ---------------------------------------------------------------
% Packages
% ---------------------------------------------------------------
\usepackage[utf8]{inputenc}
\usepackage[T1]{fontenc}
\usepackage{amsmath, amssymb}
\usepackage{graphicx}
\usepackage{booktabs}
\usepackage{array}
\usepackage{hyperref}
\usepackage{enumitem}
\usepackage{geometry}
\usepackage{caption}
\usepackage{subcaption}
\usepackage{cite}
\usepackage{url}
\usepackage{parskip}

\geometry{
  top=2.5cm,
  bottom=2.5cm,
  left=2.5cm,
  right=2.5cm
}

\hypersetup{
  colorlinks=true,
  linkcolor=blue,
  citecolor=blue,
  urlcolor=blue
}

% ---------------------------------------------------------------
% Document
% ---------------------------------------------------------------
\begin{document}

% ---------------------------------------------------------------
% Title & Author
% ---------------------------------------------------------------
\begin{center}
  {\Large \textbf{Obesity Level Prediction Using Machine Learning}}\\[1em]
  Hardikkumar Jayswal\textsuperscript{3}~\orcidlink{0000-0002-2224-5258}\\[0.5em]
  \small
  Department of Information Technology, Devang Patel Institute of Advanced Technology and Research,\\
  Charotar University of Science and Technology, Anand (Gujarat), India\\[0.3em]
  \texttt{hardikjayswal.it@charusat.ac.in}
\end{center}

% ---------------------------------------------------------------
% Abstract
% ---------------------------------------------------------------
\begin{abstract}
Obesity has become a significant global health concern due to its association with
various non-communicable diseases. Traditional methods for obesity assessment,
such as BMI, often fail to capture the complexity of the condition, highlighting the
need for more accurate predictive tools. This research utilizes machine learning
algorithms, including Random Forest, Gradient Boosting, Support Vector Machines,
and Neural Networks, in a stacking ensemble model to predict obesity levels.
Utilizing datasets from diverse populations, the model achieved a high accuracy of
96.69\%. Key features such as BMI, age, and dietary habits were identified as critical
predictors through Recursive Feature Elimination. The research findings demonstrate
the potential of advanced data-driven techniques in providing personalized insights
into obesity management and underscore the transformative role of machine learning
in public health initiatives.

\textbf{Keywords:} Machine Learning, Recursive Feature Elimination, Obesity.
\end{abstract}

% ---------------------------------------------------------------
\section{Introduction}
% ---------------------------------------------------------------

Obesity poses a major public health challenge worldwide, having escalated to epidemic
levels and affecting a vast number of individuals globally. As per the World Health
Organization, more than 650 million adults met the obesity criteria in 2016, defined by a
Body Mass Index exceeding 30~kg/m\textsuperscript{2}~\cite{who2016}. The consequences
of obesity extend beyond aesthetics, contributing substantially to the global burden of
disease, with direct links to type~2 diabetes, cardiovascular diseases, musculoskeletal
disorders, and certain types of cancer~\cite{who2016}. The global rise in obesity poses
severe health risks and imposes an enormous financial burden on healthcare systems across
both developed and developing nations~\cite{hammond2010}.

\subsection{The Global Obesity Epidemic}

The prevalence of obesity has worsened significantly in recent decades, with various
socioeconomic, environmental, and lifestyle changes fueling the increase in rates of
obesity worldwide~\cite{ng2014}. Urbanization and the rapid adoption of sedentary
lifestyles in both developed and developing nations have led to a dramatic shift in
dietary patterns. Increased consumption of energy-dense, nutrient-poor processed foods
and reduced physical activity have become common features of modern life~\cite{popkin2012}.
These lifestyle shifts have driven the rise of obesity-related diseases, such as type~2
diabetes and cardiovascular conditions, which are now some of the leading causes of
mortality globally~\cite{swinburn2011}.

In many low and middle-income countries, the obesity epidemic is particularly concerning,
as it coexists with undernutrition, creating a double burden of disease~\cite{tzioumis2014}.
The rapid spread of fast food, sugary beverages, and ultra-processed products has
exacerbated this trend, leading to an alarming increase in obesity rates, even among
children and adolescents~\cite{lobstein2015}. Despite various public health campaigns and
initiatives aimed at promoting healthier lifestyles, global efforts to curb the obesity
epidemic have had limited success~\cite{swinburn2011}. The persistence of the epidemic
underscores the need for innovative and scalable solutions that can address the root
causes of obesity and offer effective prevention and management strategies.

\subsection{Limitations of Traditional Approaches}

Traditional methods for assessing and managing obesity, such as BMI calculations
and general lifestyle assessments, have proven insufficient in addressing the complex
and multifaceted nature of the condition~\cite{hruby2015}. BMI, though widely used,
is a simplistic measure that does not account for variations in muscle mass, fat
distribution, or metabolic health. Furthermore, traditional approaches often fail to
account for individual variability in genetic predisposition, metabolic rates, and the
socio-environmental factors influencing weight gain~\cite{wells2012}. These
limitations can lead to misclassifications and delayed interventions, reducing the
overall effectiveness of obesity management strategies.

Moreover, manual assessments, which typically involve dietary surveys and physical
activity logs, are time-consuming and prone to human error~\cite{pierson2000}.
Healthcare professionals face the additional challenge of trying to tailor treatment
and prevention strategies to individual patients without the necessary data insights.
The growing complexity of obesity highlights the inadequacy of these conventional
methods, emphasizing the need for more advanced, data-driven approaches that can
provide timely and personalized interventions~\cite{burton2005}.

\subsection{Promise of Machine Learning and Big Data Analysis}

In the face of these challenges, machine learning~\cite{jayswal2023plant} and big data
analytics offer promising alternatives for improving obesity prediction, prevention,
and management. By analyzing vast and complex datasets, ML algorithms can uncover
patterns and relationships that are not immediately apparent through traditional
statistical methods~\cite{ferdowsy2021}. These advanced techniques allow for more
precise and accurate predictions of obesity risk, enabling earlier detection and
intervention~\cite{esteva2019}. Furthermore, ML-based models can handle a wide range
of variables, such as genetic information, behavioral data, and environmental
factors, to create a more holistic understanding of obesity's root causes and
progression~\cite{esteva2019}.

ML algorithms, such as decision trees, support vector machines, and neural networks,
have been shown to outperform traditional models in classifying and predicting obesity
levels based on large, diverse datasets~\cite{cervantes2020}. These algorithms can
integrate multiple types of data, including demographic information, lifestyle habits,
and medical history, allowing for the development of personalized prevention and
treatment plans tailored to individual needs~\cite{ferdowsy2021}. Additionally, big
data analysis enables healthcare providers to identify high-risk populations, monitor
trends over time, and allocate resources more effectively in the fight against
obesity~\cite{esteva2019}.

In conclusion, the integration of machine learning and big data analytics into obesity
research presents a significant opportunity to address the limitations of traditional
approaches. These technologies offer the potential for more accurate, personalized, and
scalable solutions that could have a transformative impact on global obesity prevention
and management efforts.

% ---------------------------------------------------------------
\section{Literature Review}
% ---------------------------------------------------------------

Current advances in ML have significantly enhanced the predictive modelling of obesity,
a major global health concern linked to numerous non-communicable diseases. In their
2022 study, Bimantara and Supriana explored the integration of Case-Based Reasoning
with K-Means clustering to estimate obesity levels, demonstrating the method's
efficiency by reducing computation time during the retrieval of similar cases, achieving
an accuracy of 88.27\%~\cite{bimantara2022}. Concurrently, Manna et al.\ in 2014
utilized fuzzy signatures to predict early childhood obesity risks, emphasizing the
importance of capturing complex, nonlinear interactions within data, which traditional
statistical methods might overlook~\cite{manna2014}. In their 2018 study, Pavlyshenko
investigated stacking approaches in machine learning to enhance the predictive accuracy
of obesity models, illustrating how combining multiple predictive models could lead to
more robust and accurate predictions~\cite{pavlyshenko2018}. In a similar vein,
Ferdowsy et al.\ in 2021 employed various advanced machine learning techniques,
including neural networks and decision trees, on large-scale datasets from electronic
health records and mobile devices to classify and predict obesity, demonstrating the
potential of big data in healthcare~\cite{ferdowsy2021}. Lastly, Cervantes and Palacio
in 2020 developed a logistic regression model to predict the likelihood of different BMI
categories among children in rural areas using publicly available datasets, which aids
in identifying and implementing effective local interventions~\cite{cervantes2020}.
These studies collectively highlight a shift towards leveraging complex machine learning
techniques and big data to address the multifaceted challenges of obesity, suggesting a
trend towards more personalized and precise public health strategies.

% ---------------------------------------------------------------
\section{Materials and Method}
% ---------------------------------------------------------------

\subsection{Dataset Description}

\begin{table}[ht]
\centering
\caption{Details of Dataset.}
\label{tab:dataset_details}
\begin{tabular}{llccc}
\toprule
\textbf{Dataset Name} & \textbf{Dataset Source} & \textbf{No.\ of Instances} &
\textbf{No.\ of Instances After Pre-processing} & \textbf{No.\ of Attributes} \\
\midrule
Dataset 1~\cite{palechor2019} & Study from Colombia, Peru, and Mexico & 2{,}111 & 2{,}111 & 17 \\
Dataset 2~\cite{cdc2023}      & Centers for Disease Control and Prevention & 840 & 5{,}000 & 16 \\
\bottomrule
\end{tabular}
\end{table}

\subsubsection*{a.\ Dataset 1}

This study utilizes a dataset designed to estimate obesity levels in individuals from
Colombia, Peru, and Mexico. The dataset includes responses from 2,111 participants aged
between 14 and 61, capturing variables related to diet, physical activity, and lifestyle
habits~\cite{ferdowsy2021,cervantes2020}. The study utilized a stratified sampling
approach to select participants, ensuring diverse representation across various age
brackets, genders, and socioeconomic backgrounds~\cite{cervantes2020}. The dataset is
divided into seven obesity categories, ranging from Insufficient Weight to Obesity Type
III, based on BMI calculations~\cite{who2000}.

The data were collected through a web-based survey, which included questions about
eating habits, physical activity, and personal health history. To address class
imbalance, Synthetic Minority Over-sampling Technique (SMOTE) was applied, ensuring
balanced representation across all obesity categories~\cite{ferdowsy2021,chawla2002}.
This approach enhances the model's ability to learn from minority classes, improving
prediction accuracy for underrepresented obesity levels.

The features in the dataset encompass a wide range of factors influencing obesity:

\begin{enumerate}[label=\textbf{\roman*.}]
  \item \textbf{Demographic Variables}: Gender and Age.
  \item \textbf{Anthropometric Measurements}: Height and Weight.
  \item \textbf{Family Health History}: Family History of Overweight.
  \item \textbf{Dietary Habits}: Regular intake of high-calorie foods, frequency of
        vegetable intake, number of primary meals, snacking habits, daily water intake,
        and monitoring of caloric intake.
  \item \textbf{Lifestyle Factors}: Smoking Status, Frequency of physical activity,
        duration of technology use, frequency of alcohol intake, and preferred mode of
        transportation.
\end{enumerate}

These features provide a comprehensive view of each participant's health profile,
enabling the model to capture complex interactions between different risk factors.

\begin{figure}[ht]
  \centering
  % Placeholder: Fig. 1 -- Physical Activity Frequency by Obesity Levels
  \includegraphics[width=0.45\textwidth]{fig1_physical_activity.png}
  \caption{Physical Activity Frequency by Obesity Levels.}
  \label{fig:physical_activity}
\end{figure}

\begin{figure}[ht]
  \centering
  % Placeholder: Fig. 2 -- Height vs Weight by Obesity
  \includegraphics[width=0.45\textwidth]{fig2_height_weight.png}
  \caption{Height vs Weight by Obesity.}
  \label{fig:height_weight}
\end{figure}

\begin{figure}[ht]
  \centering
  % Placeholder: Fig. 3 -- Age vs Frequency by Obesity
  \includegraphics[width=0.45\textwidth]{fig3_age_frequency.png}
  \caption{Age vs Frequency by Obesity.}
  \label{fig:age_frequency}
\end{figure}

\begin{figure}[ht]
  \centering
  % Placeholder: Fig. 4 -- Count Vs Gender Distribution
  \includegraphics[width=0.45\textwidth]{fig4_gender_distribution.png}
  \caption{Count Vs Gender Distribution.}
  \label{fig:gender_distribution}
\end{figure}

\subsubsection*{b.\ Dataset 2}

The second dataset focuses on the occurrence of obesity among kids and youths aged 2 to
19 years in the United States, containing 840 original entries and expanded to 5,840
instances after applying SMOTE~\cite{chawla2002,cdc2023}. The dataset includes
demographic factors such as age, sex, race, and Hispanic origin, allowing for an
analysis of obesity trends across different population subgroups. The data were collected
through national health surveys and physical examinations, ensuring accurate assessment
of obesity status based on BMI percentiles appropriate for children and
adolescents~\cite{cdc2023}.

By including this dataset, the study extends the predictive modeling to younger
populations, addressing the critical issue of childhood obesity. Early identification of
at-risk individuals can lead to timely interventions, reducing the likelihood of
obesity-related complications later in life~\cite{reilly2011}.

\noindent Features of the Dataset:
\begin{itemize}
  \item \textbf{Demographic Variables}: Age, Sex, Race, and Hispanic Origin.
  \item \textbf{Health Indicators}: BMI Percentile, Obesity Status (Overweight, Obese,
        Severely Obese).
  \item \textbf{Survey Data}: Dietary Intake, Physical Activity Levels, Socioeconomic
        Factors.
  \item \textbf{Year}: Data collection period (e.g., 2007--2010).
\end{itemize}

\begin{figure}[ht]
  \centering
  % Placeholder: Fig. 5 -- Trend of Average Obesity Estimate over the Years
  \includegraphics[width=0.45\textwidth]{fig5_obesity_trend.png}
  \caption{Trend of Average Obesity Estimate over the years.}
  \label{fig:obesity_trend}
\end{figure}

\begin{figure}[ht]
  \centering
  % Placeholder: Fig. 6 -- Pie Chart of Distribution of Different Factors of Obesity
  \includegraphics[width=0.45\textwidth]{fig6_pie_chart.png}
  \caption{Pie Chart of Distributing of different factor of obesity.}
  \label{fig:pie_chart}
\end{figure}

% ---------------------------------------------------------------
\section{Data Preprocessing}
% ---------------------------------------------------------------

It is one of the decisive steps in ML, as it preprocesses the raw data for modeling by
addressing missing values, transforming categorical variables, and normalizing numerical
features~\cite{kotsiantis2006}. In this study, we performed the following preprocessing
steps:

\begin{enumerate}[label=\textbf{\alph*.}]
  \item \textbf{Handling Missing Values}: Missing numerical data were imputed using the
        mean of the respective feature, while missing categorical data were filled using
        the mode. This approach preserves the dataset's size and reduces potential biases
        introduced by missing data~\cite{guyon2002}.

  \item \textbf{Addressing Class Imbalance}: Both datasets exhibited class imbalance,
        with certain obesity categories underrepresented. We applied SMOTE to
        synthetically generate new instances for minority classes, resulting in a
        balanced dataset that enhances the model's ability to learn from all
        classes~\cite{chawla2002}.

  \item \textbf{Encoding Categorical Variables}: Categorical variables were encoded
        using one-hot encoding, transforming them into binary vectors suitable for
        machine learning algorithms~\cite{pedregosa2011}. This method prevents the
        introduction of ordinal relationships where none exist.

  \item \textbf{Feature Scaling}: Numerical features were standardized using Standard
        Scaler, which scales the data to have a mean of 0 and a standard deviation of 1.
        Standardization is essential for algorithms that rely on the magnitude of
        features, such as SVM and neural networks~\cite{jain2000}.

  \item \textbf{Feature Engineering}: We created polynomial features to identify
        non-linear relationships within the data by incorporating interaction terms and
        features of higher degrees. This process enables the model to recognize more
        intricate patterns~\cite{rumelhart1986}.

  \item \textbf{Data Splitting}: The datasets were split into training and testing sets
        using an 80--20 ratio. Stratified sampling was used to maintain the class
        distribution in both sets, ensuring that each obesity category is adequately
        represented~\cite{he2009}.
\end{enumerate}

% ---------------------------------------------------------------
\section{Feature Selection}
% ---------------------------------------------------------------

Feature selection improves the performance of the model by minimizing overfitting,
enhancing accuracy, and lowering computational expenses~\cite{kohavi1997}. We utilized
Recursive Feature Elimination (RFE) alongside cross-validation to determine the most
important features~\cite{guyon2002,kohavi1997}. RFE operates by systematically
eliminating less relevant features and constructing a model based on the features that
remain. This process is repeated until the ideal number of features is achieved.

The RFE process revealed that features such as BMI, Age, Physical Activity Frequency,
Time Spent Using Technology, and Consumption of High-Calorie Foods were among the most
influential in predicting obesity levels. By focusing on these critical features, the
model becomes more interpretable and efficient.

% ---------------------------------------------------------------
\section{Model Development}
% ---------------------------------------------------------------

\textbf{Stacking Model:} In this study, a stacking ensemble learning approach was
applied to improve the accuracy of predicting obesity levels. Stacking combines multiple
classification models to optimize predictive performance by leveraging the strengths of
each model~\cite{breiman1996,friedman2001}.

\subsubsection*{a.\ Base Models}

\begin{itemize}
  \item \textbf{Random Forest Classifier (RF):} An ensemble of decision trees that
        reduces variance through averaging, mitigating overfitting~\cite{breiman1996}.
  \item \textbf{Gradient Boosting Classifier (GB):} Builds additive models in a
        forward stage-wise fashion, allowing for the optimization of arbitrary
        differentiable loss functions~\cite{friedman2001,ferdowsy2021}.
  \item \textbf{Support Vector Machine (SVM):} Effective in high-dimensional spaces and
        uses kernel functions to capture non-linear relationships~\cite{cortes1995}.
  \item \textbf{Neural Network (NN):} A multi-layer perceptron capable of modeling
        complex non-linear relationships through multiple hidden layers and activation
        functions~\cite{bishop1995,rumelhart1986}.
\end{itemize}

\subsubsection*{b.\ Meta-Model}

The outputs from the base models are combined using a Logistic Regression model, which
acts as the meta-learner~\cite{breiman1996,friedman2001}. Logistic Regression is
suitable for classification problems and provides probabilistic outputs, making it
effective for combining predictions.

\subsubsection*{c.\ Training Process}

\begin{itemize}
  \item \textbf{Cross-Validation}: The training data were split into folds, and base
        models were trained on different subsets to prevent overfitting~\cite{he2009,pedregosa2011}.
  \item \textbf{Blending Predictions}: The predictions of the base models were used as
        input features for the meta-model. This approach allows the meta-model to learn
        how to best combine the base models' outputs~\cite{breiman1996}.
\end{itemize}

\begin{figure}[ht]
  \centering
  % Placeholder: Fig. 7 -- Process Flow to predict the obesity level
  \includegraphics[width=0.85\textwidth]{fig7_process_flow.png}
  \caption{Process Flow to predict the obesity level~\cite{kumar2024,kumar2023}.}
  \label{fig:process_flow}
\end{figure}

% ---------------------------------------------------------------
\section{Performance and Evaluation}
% ---------------------------------------------------------------

\begin{table}[ht]
\centering
\caption{Model and its Performance.}
\label{tab:model_performance}
\begin{tabular}{lcc}
\toprule
\textbf{Model} & \textbf{Dataset 1} & \textbf{Dataset 2} \\
\midrule
Random Forest        & 0.93  & 0.9   \\
Gradient Boosting    & 0.94  & 0.84  \\
SVM                  & 0.93  & 0.839 \\
Neural Network       & 0.926 & 0.91  \\
Logistic Regression  & 0.87  & 0.67  \\
Stacking             & 0.969 & 0.899 \\
\bottomrule
\end{tabular}
\end{table}

The performance of the stacking model was evaluated using the testing set and
cross-validation scores. Key metrics included accuracy, precision, recall, F1-score,
and confusion matrices~\cite{sokolova2009}.

\begin{enumerate}[label=\textbf{\alph*.}]
  \item \textbf{Accuracy}: The stacking model achieved an overall accuracy of 96.69\%,
        indicating that it correctly classified obesity levels in the majority of cases.
  \item \textbf{Precision and Recall:} High precision and recall values across most
        classes suggest that the model is both accurate and sensitive in predicting
        obesity levels~\cite{sokolova2009}.
  \item \textbf{F1-Score:} The harmonic mean of precision and recall, F1-scores were
        consistently high, demonstrating balanced performance.
  \item \textbf{Confusion Matrix:} Visual analysis of confusion matrices showed minimal
        misclassifications, with most instances correctly identified~\cite{sokolova2009}.
  \item \textbf{Cross-Validation Scores:} The mean cross-validation score was 97.57\%,
        confirming the model's robustness and generalizability~\cite{he2009}.
\end{enumerate}

\subsubsection*{For dataset 1:}

\begin{table}[ht]
  \centering
  \caption{Confusion Matrix for the dataset 1.}
  \label{tab:confusion_dataset1}
  % Placeholder: The confusion matrix images for dataset 1 are presented as figures.
  % Four confusion matrices (RF, GB, SVM, NN) were included in the original paper.
  \begin{tabular}{c}
    \includegraphics[width=\textwidth]{fig_confusion_dataset1.png}
    % Uncertain content: confusion matrix image for dataset 1
  \end{tabular}
\end{table}

\subsubsection*{For dataset 2:}

\begin{table}[ht]
  \centering
  \caption{Confusion Matrix for the dataset 2.}
  \label{tab:confusion_dataset2}
  % Placeholder: The confusion matrix images for dataset 2 are presented as figures.
  \begin{tabular}{c}
    \includegraphics[width=\textwidth]{fig_confusion_dataset2.png}
    % Uncertain content: confusion matrix image for dataset 2
  \end{tabular}
\end{table}

\subsubsection*{For stacking model:}

\begin{table}[ht]
  \centering
  \caption{Confusion Matrix for the Dataset 1 and Dataset 2.}
  \label{tab:confusion_stacking}
  % Placeholder: Stacking model confusion matrices for both datasets.
  \begin{tabular}{c}
    \includegraphics[width=\textwidth]{fig_confusion_stacking.png}
    % Uncertain content: stacking confusion matrix image
  \end{tabular}
\end{table}

\subsubsection*{Comparison with Individual Models:}

The stacking model outperformed individual base models, which had lower accuracy scores
when used alone. This result highlights the effectiveness of stacking in leveraging the
strengths of different algorithms to improve predictive
performance~\cite{breiman1996,pavlyshenko2018}.

% ---------------------------------------------------------------
\section{Discussion}
% ---------------------------------------------------------------

The high performance of the stacking model demonstrates the potential of ensemble
learning methods in predicting obesity levels. By combining models that capture
different aspects of the data, the stacking approach addresses the complexity inherent
in obesity prediction~\cite{breiman1996,pavlyshenko2018}. The use of SMOTE effectively
handled class imbalance, ensuring that minority classes were adequately
represented~\cite{chawla2002}.

Feature selection through RFE contributed to model efficiency by focusing on the most
relevant features~\cite{kohavi1997}. The identified significant features align with
established risk factors for obesity, such as dietary habits, physical activity, and
sedentary behavior~\cite{tzioumis2014,lobstein2015}. This alignment enhances the
model's interpretability and applicability in real-world settings.

\textbf{Limitations:}
\begin{enumerate}[label=\textbf{\alph*.}]
  \item \textbf{Data Sources}: The datasets used may have inherent biases due to
        self-reported data and the specific populations studied~\cite{cervantes2020,cdc2023}.
  \item \textbf{Generalizability}: While the model performed well on the provided
        datasets, its applicability to other populations requires further validation.
\end{enumerate}

% ---------------------------------------------------------------
\section{Conclusion}
% ---------------------------------------------------------------

This study successfully developed a stacking ensemble model for predicting obesity
levels, achieving high accuracy and demonstrating the effectiveness of combining
multiple machine learning algorithms. The integration of SMOTE, careful data
preprocessing, and feature selection contributed to the model's performance.

The findings suggest that machine learning models can be valuable tools in public health
for early detection and intervention in obesity. Future work includes validating the
model on diverse populations, integrating additional relevant features, and exploring
the use of explainable AI techniques to enhance model transparency.

% ---------------------------------------------------------------
% References
% ---------------------------------------------------------------
\begin{thebibliography}{99}

\bibitem{who2016}
World Health Organization (2016). \url{https://www.who.int/news-room/fact-sheets/detail/obesity-and-overweight}.

\bibitem{hammond2010}
Hammond, R.~A., \& Levine, R.: The economic impact of obesity in the United States.
\textit{Diabetes, Metabolic Syndrome and Obesity: Targets and Therapy}, 2010, 3, 285--295.

\bibitem{ng2014}
Ng, M., Fleming, T., Robinson, M., et al.: Global, regional, and national prevalence of
overweight and obesity in children and adults during 1980--2013: A systematic analysis.
\textit{The Lancet}, 2014, 384(9945), 766--781.

\bibitem{popkin2012}
Popkin, B.~M., Adair, L.~S., \& Ng, S.~W.: Global nutrition transition and the pandemic
of obesity in developing countries. \textit{Nutrition Reviews}, 2012, 70(1), 3--21.

\bibitem{swinburn2011}
Swinburn, B.~A., et al.: The global obesity pandemic: Shaped by global drivers and local
environments. \textit{The Lancet}, 2011, 378(9793), 804--814.

\bibitem{tzioumis2014}
Tzioumis, E., \& Adair, L.~S.: Childhood dual burden of under- and overnutrition in
low- and middle-income countries: A critical review. \textit{Food and Nutrition Bulletin},
2014, 35(2), 230--243.

\bibitem{lobstein2015}
Lobstein, T., Jackson-Leach, R., Moodie, M.~L., et al.: Child and adolescent obesity:
Part of a bigger picture. \textit{The Lancet}, 2015, 385(9986), 2510--2520.

\bibitem{hruby2015}
Hruby, A., \& Hu, F.~B.: The epidemiology of obesity: A big picture.
\textit{Pharmacoeconomics}, 2015, 33(7), 673--689.

\bibitem{wells2012}
Wells, J.~C.~K., et al.: Beyond the BMI: Prospects for the use of body composition
indices in the assessment of health risk. \textit{Obesity Reviews}, 2012, 13(s2),
108--120.

\bibitem{pierson2000}
Pierson, R.~N., Jr., \& Wang, J.: The assessment of body composition in severe obesity.
\textit{Obesity Surgery}, 2000, 10(2), 103--104.

\bibitem{burton2005}
Burton, N.~W., et al.: Reliability and validity of a modified version of the Active
Australia physical activity survey. \textit{Annals of Leisure Research}, 2005, 8(3--4),
205--218.

\bibitem{ferdowsy2021}
Ferdowsy, F., Rahman, M.~A., \& Adnan, M.~N.: A machine learning approach for obesity
risk prediction. \textit{Current Research in Behavioral Sciences}, 2021, 2, 100053.

\bibitem{esteva2019}
Esteva, A., Robicquet, A., Ramsundar, B., et al.: A guide to deep learning in
healthcare. \textit{Nature Medicine}, 2019, 25(1), 24--29.

\bibitem{cervantes2020}
Cervantes, R.~C., \& Palacio, U.~M.: Estimation of obesity levels based on
computational intelligence. \textit{Informatics in Medicine Unlocked}, 2020, 21, 100472.

\bibitem{bimantara2022}
Bimantara, I.~M.~S., \& Supriana, I.~W.: Case-Based Reasoning (CBR) for Obesity Level
Estimation Using K-Means Indexing Method. \textit{Jurnal Ilmiah Kursor}, 2022, 11(4),
145--156.

\bibitem{manna2014}
Manna, S., Rana, S., Saha, G., \& Singh, P.~K.: Understanding early childhood obesity
risks using fuzzy signatures. In \textit{2014 IEEE International Conference on Fuzzy
Systems}, 2014, pp.\ 2480--2487, IEEE.

\bibitem{pavlyshenko2018}
Pavlyshenko, B.~M.: Using stacking approaches for machine learning models. In
\textit{2018 IEEE Second International Conference on Data Stream Mining \& Processing
(DSMP)}, 2018, pp.\ 255--258, IEEE.

\bibitem{who2000}
World Health Organization, 2000, \textit{Obesity: Preventing and Managing the Global
Epidemic} (No.\ 894).

\bibitem{chawla2002}
Chawla, N.~V., Bowyer, K.~W., Hall, L.~O., \& Kegelmeyer, W.~P.: SMOTE: Synthetic
minority over-sampling technique. \textit{Journal of Artificial Intelligence Research},
2002, 16, 321--357.

\bibitem{cdc2023}
Centers for Disease Control and Prevention (CDC), 2023, Obesity among children and
adolescents aged 2--19 years: United States. Retrieved from CDC website.

\bibitem{reilly2011}
Reilly, J.~J., \& Kelly, J.: Long-term impact of overweight and obesity in childhood and
adolescence on morbidity and premature mortality in adulthood: Systematic review.
\textit{International Journal of Obesity}, 2011, 35(7), 891--898.

\bibitem{kotsiantis2006}
Kotsiantis, S., Kanellopoulos, D., \& Pintelas, P.: Data preprocessing for supervised
learning. \textit{International Journal of Computer Science}, 2006, 1(2), 111--117.

\bibitem{guyon2002}
Guyon, I., Weston, J., Barnhill, S., \& Vapnik, V.: Gene selection for cancer
classification using support vector machines. \textit{Machine Learning}, 2002,
46(1--3), 389--422.

\bibitem{pedregosa2011}
Pedregosa, F., Varoquaux, G., Gramfort, A., et al.: Scikit-learn: Machine learning in
Python. \textit{Journal of Machine Learning Research}, 2011, 12, 2825--2830.

\bibitem{jain2000}
Jain, A.~K., Duin, R.~P., \& Mao, J.: Statistical pattern recognition: A review.
\textit{IEEE Transactions on Pattern Analysis and Machine Intelligence}, 2000, 22(1),
4--37.

\bibitem{rumelhart1986}
Rumelhart, D.~E., Hinton, G.~E., \& Williams, R.~J.: Learning representations by
back-propagating errors. \textit{Nature}, 1986, 323(6088), 533--536.

\bibitem{he2009}
He, H., \& Garcia, E.~A.: Learning from imbalanced data. \textit{IEEE Transactions on
Knowledge and Data Engineering}, 2009, 21(9), 1263--1284.

\bibitem{kohavi1997}
Kohavi, R., \& John, G.~H.: Wrappers for feature subset selection.
\textit{Artificial Intelligence}, 1997, 97(1--2), 273--324.

\bibitem{breiman1996}
Breiman, L.: Stacked regressions. \textit{Machine Learning}, 1996, 24(1), 49--64.

\bibitem{friedman2001}
Friedman, J.~H.: Greedy function approximation: A gradient boosting machine.
\textit{Annals of Statistics}, 2001, 29(5), 1189--1232.

\bibitem{cortes1995}
Cortes, C., \& Vapnik, V.: Support-vector networks. \textit{Machine Learning}, 1995,
20(3), 273--297.

\bibitem{bishop1995}
Bishop, C.~M.: \textit{Neural Networks for Pattern Recognition}. Oxford University
Press, 1995.

\bibitem{sokolova2009}
Sokolova, M., \& Lapalme, G.: A systematic analysis of performance measures for
classification tasks. \textit{Information Processing \& Management}, 2009, 45(4),
427--437.

\bibitem{palechor2019}
Palechor, F.~M., \& de la Hoz Manotas, A.: Dataset for estimation of obesity levels
based on eating habits and physical condition in individuals from Colombia, Peru and
Mexico. \textit{Data in Brief, 25}, 104344.
\url{https://doi.org/10.1016/j.dib.2019.104344}

\bibitem{kumar2024}
Kumar, A., Radhika, R., \& Kumar, K.: Obesity Prediction Using Machine Learning
Technique. In \textit{2024 2nd International Conference on Networking and
Communications (ICNWC)}, 2024, pp.\ 1--7. IEEE.

\bibitem{kumar2023}
Kumar, Y., Kaur, I., Odeh, M., Bhargav, P., \& Ijaz, M.~F.: Detection and Diagnosis of
Different Types of Obesity Using Machine Learning-Based Approaches. In \textit{2023 2nd
International Engineering Conference on Electrical, Energy, and Artificial Intelligence
(EICEEAI)}, 2023, pp.\ 1--7. IEEE.

\bibitem{jayswal2023plant}
Jayswal, H.~S., \& Chaudhari, J.~P.: Plant Diseases Detection and Classification Using
Machine Learning, Deep Learning, Spectroscopy. In: Senjyu, T., So-In, C., Joshi, A.
(eds) \textit{Smart Trends in Computing and Communications. SmartCom 2023. Lecture
Notes in Networks and Systems}, vol.\ 650. Springer, Singapore.
\url{https://doi.org/10.1007/978-981-99-0838-7_51}

\bibitem{jayswal2023leaf}
Jayswal, H.~S., \& Chaudhari, J.~P.: Plant Leaf Diseases Detection and Classification
Using Spectroscopy. In: So-In, C., Londhe, N.~D., Bhatt, N., Kitsing, M. (eds)
\textit{Information Systems for Intelligent Systems. Smart Innovation, Systems and
Technologies}, vol.\ 324. Springer, Singapore.
\url{https://doi.org/10.1007/978-981-19-7447-2_42}

\end{thebibliography}

\end{document}